\subsection{The Euler-factor criterion: which rows have a $p$-adic limit, and
what it is}\label{sec:eulercrit}

\S\ref{sec:conjDproof} identified the $3$-adic limits of $\mathbf B,\mathbf C,
\mathbf F$ under the hypothesis that every $V_d$ occurring in the source is
prime to $p$ --- automatic at $p=3$ because $\chim3(3)=0$.  This subsection
removes that hypothesis and replaces it by a criterion that is visible on the
source table, Table~\ref{tab:sources}, alone.  Full details, the exact
verification record and the negative results are in
\texttt{consolidation/EULER\_CRITERION.md}; scripts in
\texttt{lattice/euler\_criterion/}.

\subsubsection*{The two placements, and the depleted Eisenstein series}

Let the source be $\Phi=P(V)E$ with
\begin{equation}\label{eq:Epsiphi}
E=E^{\psi,\varphi}_{w+2}
=\delta(\psi)L(-w-1,\varphi)+\sum_{n\ge1}\Bigl(\sum_{d\mid n}\psi(n/d)\varphi(d)d^{\,w+1}\Bigr)q^n,
\qquad L(E,s)=L(\psi,s)L(\varphi,s-w-1),
\end{equation}
$\delta(\psi)=1$ iff $\psi=\mathbf 1$, and $P(V)=\sum_dc_dV_d$ with Mellin
polynomial $P(s)=\sum_dc_dd^{-s}$.  Inner placement
$\sigma^{(w+1)}_\chi$ in Table~\ref{tab:sources} means $(\psi,\varphi)=(\chi,\mathbf 1)$;
outer placement $\widetilde\sigma^{(w+1)}_\chi$ means $(\psi,\varphi)=(\mathbf 1,\chi)$.
Applying $\thq^{-(w+1)}$ transfers the exponent to the complementary divisor:
\begin{equation}\label{eq:Ecal}
\mathcal E:=\thq^{-(w+1)}E=\sum_{n\ge1}\Bigl(\sum_{e\mid n}\psi(e)e^{-(w+1)}\varphi(n/e)\Bigr)q^n ,
\qquad \Theta=\thq^{-(w+1)}\Phi=P\bigl(V_\bullet/\bullet^{\,w+1}\bigr)\mathcal E .
\end{equation}
So $\psi$ --- the character of the \emph{unshifted} $L$-factor --- is the one
carried by the variable with the negative exponent, and it is that variable
which the $p$-adic theory depletes:
\begin{equation}\label{eq:depl}
\mathcal E^{[p]}:=\Bigl(1-\psi(p)p^{-(w+1)}V_p\Bigr)\mathcal E
=\sum_{n\ge1}\Bigl(\sum_{e\mid n,\ p\nmid e}\psi(e)e^{-(w+1)}\varphi(n/e)\Bigr)q^n .
\end{equation}
Write $\mathcal E_p(s):=1-\psi(p)p^{-s}$ for the $p$-Euler factor of $L(\psi,s)$;
under $p^{-s}\mapsto V_p/p^{w+1}$ it is exactly the operator in \eqref{eq:depl}.

\begin{lemma}[overconvergent avatar, general placement]\label{lem:avatargen}
\Lit{} Decompose $\psi=\psi^{(p)}\psi_p$ into its prime-to-$p$ and $p$-power
parts.  Then $\mathcal E^{[p]}$ is the $q$-part of the Coleman--Mazur Eisenstein
family with tame characters $(\varphi,\psi^{(p)})$ at the weight-character
$\kappa_0=\psi_p\omega^{-w}\langle\cdot\rangle^{-w}\in\mathcal W^+$, the point
``classical weight $-w$''; in particular $G^{(p)}:=\kappa_p+\mathcal E^{[p]}$ is
overconvergent of that weight, with
\begin{equation}\label{eq:kappa}
\kappa_p=\begin{cases}\dfrac12\,L_p\bigl(w+1,\ \psi\omega^{-w}\bigr), & \varphi=\mathbf 1,\\[6pt]
0, & \varphi\neq\mathbf 1 .\end{cases}
\end{equation}
\end{lemma}
\begin{proof}
$\psi(e)e^{-(w+1)}=\psi^{(p)}(e)\cdot[\psi_p(e)e^{-w}]\cdot e^{-1}$ on
$e\in\Z_p^\times$ identifies the family and the weight
\cite{Coleman1997,ColemanMazur1998}.  The constant term of the classical
member of weight $k$ is $\delta(\varphi)L(1-k,\psi)$, $p$-stabilised to
$-\tfrac12(1-\psi(p)p^{k-1})B_{k,\psi}/k=\tfrac12L_p(1-k,\psi\omega^k)$
\cite[Thm.~5.11]{Washington1997}; it vanishes identically in $k$ when
$\varphi\neq\mathbf 1$, and continuity on $\mathcal W^+$ gives \eqref{eq:kappa}
at $k=-w$.  When $\varphi=\mathbf 1$ the parity condition
$\psi\varphi(-1)=(-1)^w$ of \eqref{eq:Epsiphi} forces $\psi\omega^{-w}$ to be
\emph{even}, so $L_p(\cdot,\psi\omega^{-w})\not\equiv0$.
\end{proof}

\begin{lemma}[$V_p$ preserves overconvergence]\label{lem:Vp}
\Lit{} For $0<v<p/(p+1)$ the Frobenius $V_p$ maps $M^\dagger_k(v)$ to
$M^\dagger_k(v/p)$ \cite{Coleman1996}.  Hence any polynomial in the $V_d$,
$d$ arbitrary, carries an overconvergent form to an overconvergent form,
at the cost of a smaller strict neighbourhood of the ordinary locus and a
level multiplied by a power of $p$.
\end{lemma}

\subsubsection*{The criterion}

\begin{namedthm}[F: Euler-factor criterion]\label{thm:F}
Let $(\Gamma,t,F)$ be a modular Ap\'ery row of weight $w$ and level $N$ with
Eisenstein source $\Phi=P(V)E^{\psi,\varphi}_{w+2}$, and fix a prime $p$.
Assume
\begin{enumerate}[label=\textup{(\alph*)}]
\item $\mathcal E_p(s)\mid P(s)$ in $\Q[p^{-s}]$, i.e.\ $P$ vanishes at
$p^{-s}=\psi(p)^{-1}$ \textup{(}vacuous when $p\mid\operatorname{cond}\psi$\textup{)};
put $Q:=P/\mathcal E_p$;
\item $t$ is a Hauptmodul of a genus-zero $\Gamma_t\supseteq\Gamma_0(N')$,
$N'$ the level of $Q(V_\bullet/\bullet^{w+1})G^{(p)}$, and
$H_{\xi^*}=F(\Theta-\xi^*)$ is $\Gamma_t$-invariant;
\item $t\,j\in\Z_{(p)}[[t]]$ with unit constant term;
\item $v_p(a_n)=O(\log n)$ and $\sigma_p>0$.
\end{enumerate}
Then $b_n/a_n$ converges in $\Q_p$ and
\begin{equation}\label{eq:xivalue}
\boxed{\ \xi_p=-\,Q(w+1)\,\kappa_p\ }
\end{equation}
with $\kappa_p$ as in \eqref{eq:kappa}.  If \textup{(a)} fails, no $\xi$ makes
$H_\xi$ overconvergent.  \Proved{} modulo \textup{(b)}, \textup{(d)} and the
citations of Lemmas~\ref{lem:avatargen}--\ref{lem:Vp} and
\cite[4.4.2]{Katz1973}.
\end{namedthm}

\begin{proof}[Proof sketch]
By \eqref{eq:Ecal}--\eqref{eq:depl}, (a) says precisely
$\Theta=Q(V_\bullet/\bullet^{w+1})\mathcal E^{[p]}
 =Q(V_\bullet/\bullet^{w+1})G^{(p)}-Q(w+1)\kappa_p$, a \emph{finite} oldform
combination of an overconvergent form (Lemmas~\ref{lem:avatargen},
\ref{lem:Vp}); note $Q(V_\bullet/\bullet^{w+1})$ acts on constants by
$Q(w+1)$.  Then $H_{\xi^*}$ with $\xi^*=-Q(w+1)\kappa_p$ is overconvergent of
weight $0$, hence a rigid function on a strict neighbourhood of $D_\infty$;
(c) puts $\{|t|_p\le1\}$ inside $D_\infty$ and (b) makes the $t$-expansion
$B-\xi^*A$ converge on $|t|_p\le\rho$ for some $\rho>1$.  By
Lemma~\ref{lem:overconv} and (d), $\xi^*=\xi_p$.  Conversely for
$\xi\ne\xi^*$, $H_\xi-H_{\xi^*}=(\xi^*-\xi)F$ would be overconvergent of both
weight $0$ and weight $w$, contradicting \cite[4.4.2]{Katz1973}.  If (a)
fails, $P/\mathcal E_p$ is the infinite series
$\sum_k\psi(p)^kp^{-k(w+1)}V_p^k$, whose partial sums leave every fixed strict
neighbourhood of the ordinary locus.
\end{proof}

\begin{proposition}[cuspidal clause]\label{prop:cusp}
\Proved{} Let instead $\Phi=P(V)f$ with $f$ a normalised cuspidal eigenform of
weight $w+2$ and $U_pf=a_pf$, and put $Y=V_p/p^{w+1}$.  If $(1-a_pY)\mid P(Y)$
\textup{(}automatic when $p^2\mid N$, so $a_p=0$\textup{)}, then
$\Theta=Q(Y)\thq^{-(w+1)}f^{[p]}$ with $f^{[p]}=f-a_pV_pf$ is overconvergent
of weight $-w$ \emph{and has no constant term}, by Coleman's
$\theta^{\,w+1}$-surjectivity onto $U_p$-depleted overconvergent forms
\cite{Coleman1996}.  Hence, under \textup{(b)}--\textup{(d)},
\[
\xi_p=0 :
\]
the $p$-adic Ap\'ery limit of a cuspidal row vanishes whenever it exists.
\end{proposition}

\subsubsection*{Consequence: Conjecture~D, with an Euler correction}

\begin{corollary}[$p$-adic rational factor]\label{cor:rp}
\Proved{} If $\varphi=\mathbf 1$ then, with $\Lambda:=L(\psi,w+1)$ and
$\Lambda_p:=L_p(w+1,\psi\omega^{-w})$,
\[
\xi_\infty=L(\Phi,w+1)=-\tfrac12P(w+1)\Lambda=:r_\infty\Lambda,\qquad
\xi_p=-\tfrac12Q(w+1)\Lambda_p=:r_p\Lambda_p,\qquad
\boxed{\ r_p=\frac{r_\infty}{\mathcal E_p(w+1)}\ }
\]
The correction $\mathcal E_p(w+1)=1-\psi(p)p^{-(w+1)}$ depends only on
$(\psi,p,w)$, not on the row.  Hence any two rows with the same Eisenstein
series $E$ satisfying \textup{(a)}--\textup{(d)} at $p$ have
$r\,\xi'_p=r'\,\xi_p$: Conjecture~\ref{conj:D} holds for them, with
$\Lambda_p$ the $p$-adic avatar of $\Lambda$.  If $\varphi\neq\mathbf 1$ then
$\xi_p=0$.
\end{corollary}

\begin{remark}[what the criterion sees on Table~\ref{tab:sources}]
\VerifiedR{exact, all twelve families, $p=2,3,5$}
Divisibility (a) holds \emph{exactly} at the primes dividing $c$:
\[
P_{\mathbf A}=1-X,\quad P_{\mathbf F}=(1+X)(1-8X),\quad
P_\alpha=(1-X)(1-16X)(1-9Y),
\]
\[
P_\varepsilon=(1-X)(1-4X)(1-16X),\qquad P_\delta=(1-Y)(1-14X+16X^2),
\]
$(X=2^{-s},Y=3^{-s})$, while $P_{\mathbf C},P_{\mathbf B}$ do not vanish at
$X=-1=\chim3(2)^{-1}$, $P_\eta$ does not vanish at $X=-1=\chi_5(2)^{-1}$, and
$P_\gamma$ vanishes at neither $X=1$ nor $Y=1$.  Thirty-six tests, no
exception: \emph{$\mathcal E_p\mid P\iff p\mid c$}.  This is an algebraic
reading of the slope set of Proposition~\ref{prop:C}.
\end{remark}

\begin{verification}[the census]\label{ver:eulercensus}
\VerifiedR{$\ge p^{2991}$ for every row; \texttt{lattice/euler\_criterion/}}
$\kappa_p$ is computed from Washington's Thm.~5.11 series at integer $s$
\textup{(}so all binomial coefficients are exact integers\textup{)} to
precision $p^{3060}$, validated against the interpolation property at
$n=1,\dots,9$ for each character used, and cross-checked independently by the
Kummer sequence of exact generalised Bernoulli numbers
$-(1-\psi(p)p^{k-1})B_{k,\psi}/(2k)$, $k\equiv-w$, whose $p$-adic distance to
$\kappa_p$ decreases by exactly one power of $p$ per step.  Rows are generated
exactly over $\Q$ to $N$ with $\sigma_pN\ge3040$.  Every entry below is
agreement to the full precision available, not a measured discrepancy.
Here $\zeta_p(s)=L_p(s,\mathbf 1)$; $\zeta_2(2)$ is Calegari's
$L_2(2,\chim4)$ in Washington's normalisation.

\smallskip
\begin{center}\small
\renewcommand{\arraystretch}{1.2}
\begin{tabular}{@{}l c c c c c c c@{}}
\toprule
row & $p$ & $(\psi,\varphi)$ & $\mathcal E_p(w{+}1)$ & $Q(w{+}1)$ & $\xi_p$ & $N$ & digits\\
\midrule
$\mathbf A$ $(7,2,-8)$    & 2 & $(\mathbf 1,\chim3)$ & $3/4$ & $1$ & $0$ & $1014$ & $3003$\\
$\mathbf E$ $(12,4,32)$   & 2 & $(\chim4,\mathbf 1)$ & $1$ & $-1$ & $\tfrac12\zeta_2(2)$ & $608$ & $3026$\\
$\mathbf F$ $(17,6,72)$   & 2 & $(\chim3,\mathbf 1)$ & $5/4$ & $-1$ & $\tfrac12L_2(2,\chi_{12})$ & $1014$ & $3005$\\
$\mathbf F$ $(17,6,72)$   & 3 & $(\chim3,\mathbf 1)$ & $1$ & $-5/4$ & $\tfrac58\zeta_3(2)$ & $1520$ & $3025$\\
$\mathbf B$ $(9,3,27)$    & 3 & $(\chim3,\mathbf 1)$ & $1$ & $-1$ & $\tfrac12\zeta_3(2)$ & $1014$ & $3029$\\
$\mathbf C$ $(10,3,9)$    & 3 & $(\chim3,\mathbf 1)$ & $1$ & $-1$ & $\tfrac12\zeta_3(2)$ & $1520$ & $3025$\\
$\alpha$ Domb $(10,4,64)$ & 2 & $(\mathbf 1,\mathbf 1)$ & $7/8$ & $-2/3$ & $\tfrac13\zeta_2(3)$ & $507$ & $2991$\\
$\varepsilon$ $(12,4,16)$ & 2 & $(\mathbf 1,\mathbf 1)$ & $7/8$ & $-1/2$ & $\tfrac14\zeta_2(3)$ & $760$ & $3000$\\
$\delta$ $(7,3,81)$       & 3 & $(\mathbf 1,\mathbf 1)$ & $26/27$ & $-1/2$ & $\tfrac14\zeta_3(3)$ & $760$ & $3030$\\
$\zeta$ $(9,3,-27)$       & 3 & $(\chim3,\chim3)$ & $1$ & $1$ & $0$ & $1014$ & $3023$\\
$\eta$ $(11,5,125)$       & 5 & $(\chi_5,\mathbf 1)$ & $1$ & $-1$ & $\tfrac12\zeta_5(3)$ & $1014$ & $3027$\\
\midrule
Cooper $s_{18}$           & 3 & --- & $1$ & $-1$ & $\tfrac12\zeta_3(2)$ & $3040$ & $3025$\\
Zudilin (Catalan)         & 2 & --- & $1$ & $-2$ & $\zeta_2(2)$ & $380$ & $3015$\\
cusp row $L(f,2)$, $N{=}12$ & 2 & cuspidal, $a_2{=}0$ & $1$ & $1$ & $0$ & $1520$ & $3027$\\
\bottomrule
\end{tabular}
\end{center}
\smallskip

Conversely, where (a) fails there is no $p$-adic limit: the increments
$v_p(b_n/a_n-b_{n-1}/a_{n-1})$ stay bounded and negative for $\mathbf B,
\mathbf C$ at $p=2$, $\gamma$ at $p=2,3$, $\delta$ and $\eta$ at $p=2$,
$\alpha$ at $p=3$, the cusp row at $p=3$, and for the Domb cuspidal apparatus
$\bigl(f_*=f_6-4V_2f_6$, $a_2=-2$, $a_3=-3$, $P=1-4V_2\bigr)$ at both $p=2$
and $p=3$, where $1-4V_2$ is divisible by neither $1+2Y$ nor $1+3Y$.  Nor do
$p$-power towers $n=ap^s$ rescue the failed cases: for $\mathbf C$ at $p=2$
and $\alpha$ at $p=3$ the tower increments \emph{decrease} linearly in $s$.
\end{verification}

\begin{remark}[reading of the table]
\begin{enumerate}[leftmargin=1.6em]
\item $\alpha$ and $\varepsilon$ at $p=2$ carry the same $\Lambda_2=\zeta_2(3)$
with $r_2=\tfrac13,\tfrac14$, ratio $\tfrac43$ --- the archimedean ratio
$(7/24)/(7/32)$, both factors divided by the same $\mathcal E_2(3)=\tfrac78$.
This settles the $p\mid d$ case left open at the end of \S\ref{sec:conjDproof}.
\item $\mathbf F$ at $p=2$ is the first case in which the Euler factor is
genuinely cancelled: $\chim3(2)=-1$, $\mathcal E_2=1+X$, and $Q=1-8V_2$ still
contains $V_p$ --- harmless by Lemma~\ref{lem:Vp}.  The twist by
$\omega^{-w}=\omega=\chim4$ is visible in the answer, $\chi_{12}=\chim3\chim4$.
\item Outer placement kills the $p$-adic period: $\mathbf A$ has
$\xi_\infty=\zeta(2)/4$ but $\xi_2=0$, and $\zeta$ has $\xi_3=0$.  Both are
$\varphi\neq\mathbf 1$; equivalently $\Lambda_p=L_p(\cdot,\psi\omega^{-w})$ is
the $L_p$ of an \emph{odd} character, identically zero.
\item Zudilin's Catalan row is not in Table~\ref{tab:sources} --- it is
hypergeometric, non-integral ($\kappa_2=4$, $c=1$, $\sigma_2=8$) --- yet the
value formula applies with $Q(2)=-2$, giving $\xi_2=\zeta_2(2)=2\xi_2^{\mathbf E}$,
matching the archimedean ratio $G:\tfrac12G$.  This is a new
Conjecture~\ref{conj:D} pair at $p=2$.
\item Cooper's $s_{18}$ joins $\mathbf B,\mathbf C,\mathbf F$ on
$\Lambda_3=\zeta_3(2)$ with $r_3=\tfrac12$, upgrading the ``aligned'' entry of
the slope census to an identified value.
\end{enumerate}
\end{remark}
